\documentclass[11pt,a4paper]{article}
\usepackage[utf8]{inputenc}
\usepackage[T1]{fontenc}
\usepackage{geometry}
\usepackage{hyperref}
\usepackage{tabularx}
\usepackage{booktabs}
\usepackage{enumitem}
\geometry{margin=1in}
\hypersetup{colorlinks=true, linkcolor=blue, urlcolor=blue}

\title{Solar System Project: Functional Explanation and Criteria Fulfillment}
\author{Project Documentation}
\date{\today}

\begin{document}
\maketitle
\tableofcontents
\newpage

\section{Project Functional Overview}

This project is a real-time 3D Solar System simulation built with C++ and Modern OpenGL (Core Profile 3.3). The program models the Sun, planets, moons, orbit paths, and helper visuals (axes and grid), then renders them with lighting, textures, and camera interaction.

\subsection{Main Runtime Flow}
\begin{enumerate}[leftmargin=*]
    \item Initialize GLFW window, OpenGL context, callbacks, and global render state in \texttt{src/main.cpp}.
    \item Load shaders, create meshes, load textures, and build scene objects in \texttt{src/solar\_system.cpp}.
    \item Run frame loop:
    \begin{itemize}[leftmargin=*]
        \item Compute \texttt{deltaTime}
        \item Process input (camera and feature toggles)
        \item Update simulation state (planet rotation/revolution)
        \item Render skybox, sun, planets, rings, orbits, axes, and grid
    \end{itemize}
    \item Exit and release resources.
\end{enumerate}

\subsection{Technical Architecture}
\begin{itemize}[leftmargin=*]
    \item \textbf{Entry and loop}: \texttt{src/main.cpp}
    \item \textbf{Scene manager}: \texttt{src/solar\_system.cpp}
    \item \textbf{Planet physics and transforms}: \texttt{src/planet.cpp}
    \item \textbf{Orbit/axis/grid generation and rendering}: \texttt{src/orbit\_renderer.cpp}
    \item \textbf{Geometry generation}: \texttt{src/sphere.cpp}, \texttt{src/ring.cpp}
    \item \textbf{Classic algorithms}: \texttt{src/algorithms/dda.cpp}, \texttt{src/algorithms/bresenham.cpp}, \texttt{src/algorithms/midpoint\_circle.cpp}
    \item \textbf{Shaders}: \texttt{shaders/*.vert}, \texttt{shaders/*.frag}
\end{itemize}

\section{Functional Explanation by Module}

\subsection{1) Input, Window, and Frame Loop}
\texttt{src/main.cpp} is responsible for:
\begin{itemize}[leftmargin=*]
    \item Creating the application window and OpenGL context (GLFW).
    \item Registering input callbacks (mouse, keyboard, scroll, resize).
    \item Computing frame time and ensuring smooth update/render behavior.
    \item Calling \texttt{solarSystem.update(deltaTime)} and \texttt{solarSystem.render(...)} each frame.
\end{itemize}

\subsection{2) Scene Initialization and Rendering Coordination}
\texttt{src/solar\_system.cpp} is the central orchestrator:
\begin{itemize}[leftmargin=*]
    \item Loads all shaders and creates shared geometry objects.
    \item Builds planets, moons, rings, and orbit render data.
    \item Applies runtime toggles: pause, orbit visibility, axes/grid visibility, transformation demo.
    \item Coordinates render order: skybox $\rightarrow$ sun $\rightarrow$ planets/rings $\rightarrow$ helper visuals.
\end{itemize}

\subsection{3) Planet State, Motion, and Transformations}
\texttt{src/planet.cpp} handles:
\begin{itemize}[leftmargin=*]
    \item Planet properties (radius, orbital radius, spin speed, revolution speed, axial tilt).
    \item Time-based animation updates for spin and revolution.
    \item Model matrix construction using translation, rotation, and scaling.
    \item Additional transformation demo functions: reflection and shear.
\end{itemize}

\subsection{4) Orbit and Helper Primitive Rendering}
\texttt{src/orbit\_renderer.cpp} provides:
\begin{itemize}[leftmargin=*]
    \item Orbit creation from Midpoint Circle output.
    \item Axis creation from DDA output.
    \item Grid creation from Bresenham output.
    \item Drawing with OpenGL primitives (line loops and points).
\end{itemize}

\subsection{5) Mesh and Visual Components}
\begin{itemize}[leftmargin=*]
    \item \texttt{src/sphere.cpp}: procedural UV sphere mesh for planets/sun.
    \item \texttt{src/ring.cpp}: ring mesh for Saturn-like ring rendering.
    \item \texttt{src/skybox.cpp}: star background rendering.
    \item \texttt{src/shader.cpp}: shader loading, compilation, linking, and uniform utilities.
\end{itemize}

\section{Criteria Fulfillment (With Where It Is Implemented)}

\begin{tabularx}{\textwidth}{@{}p{0.25\textwidth}p{0.12\textwidth}p{0.63\textwidth}@{}}
\toprule
\textbf{Required Criterion} & \textbf{Status} & \textbf{Where Fulfilled in Project} \\
\midrule
Use of basic graphics primitives (points, lines, polygons, circles) & Fulfilled &
Points: \texttt{src/orbit\_renderer.cpp} with \texttt{GL\_POINTS} for axes/grid. \\
& &
Lines/Circles: \texttt{GL\_LINE\_LOOP} for circular orbits in \texttt{src/orbit\_renderer.cpp}; orbit points generated by Midpoint Circle in \texttt{src/algorithms/midpoint\_circle.cpp}. \\
& &
Polygons: triangle meshes via \texttt{glDrawElements(GL\_TRIANGLES,...)} in \texttt{src/sphere.cpp} and \texttt{src/ring.cpp}. \\
\midrule
Implementation of at least two graphics algorithms (DDA, Bresenham, Midpoint Circle) & Fulfilled (3/3) &
DDA implementation: \texttt{src/algorithms/dda.cpp}, used by axis generation in \texttt{OrbitRenderer::generateAxes}. \\
& &
Bresenham implementation: \texttt{src/algorithms/bresenham.cpp}, used by grid generation in \texttt{OrbitRenderer::generateGrid}. \\
& &
Midpoint Circle implementation: \texttt{src/algorithms/midpoint\_circle.cpp}, used by \texttt{OrbitRenderer::addOrbit}. \\
\midrule
Use of 2D transformations (translation, rotation, scaling, reflection, shear) & Fulfilled &
Translation/Rotation/Scaling are composed in planet model matrices inside \texttt{src/planet.cpp} (orbit placement, axial tilt/spin, radius scaling). \\
& &
Reflection and shear are provided in \texttt{Planet::getReflectedModelMatrix} and \texttt{Planet::getShearedModelMatrix}, and visualized in transform demo mode from \texttt{src/solar\_system.cpp}. \\
\midrule
Inclusion of at least one moving/animated object & Fulfilled &
Continuous animation is updated per-frame using \texttt{deltaTime} in \texttt{src/main.cpp} and \texttt{SolarSystem::update}. Planets and moons continuously revolve/spin in \texttt{src/planet.cpp}. \\
\midrule
Meaningful real-world themes (scenery, monument, structure, or creative digital art) & Fulfilled &
Theme is a scientifically meaningful real-world structure: the Solar System (Sun, planets, moons, orbits, textured bodies, sky background), implemented across \texttt{src/solar\_system.cpp}, \texttt{src/planet.cpp}, \texttt{src/skybox.cpp}. \\
\midrule
Project must be original and significantly different from weekly lab tasks & Fulfilled (justified) &
The project integrates multiple advanced components beyond a typical isolated lab: full real-time scene management, camera system, lighting and textures, hierarchical motion (planet-moon), multiple algorithms in one pipeline, interactive controls, and transformation demo features. \\
\bottomrule
\end{tabularx}

\section{Technology Overview}

\subsection{Libraries and Tools}
\begin{itemize}[leftmargin=*]
    \item \textbf{C++17}: core implementation language.
    \item \textbf{OpenGL 3.3 Core}: rendering API.
    \item \textbf{GLFW}: window/context/input management.
    \item \textbf{GLAD}: OpenGL function loader.
    \item \textbf{GLM}: vector/matrix math and transformations.
    \item \textbf{stb\_image}: texture image loading.
    \item \textbf{CMake}: cross-platform build system.
\end{itemize}

\subsection{Rendering Stack}
\begin{itemize}[leftmargin=*]
    \item Vertex and fragment shaders for planets, sun, rings, skybox, and lines.
    \item VAO/VBO/EBO pipeline for mesh and helper geometry.
    \item Depth testing, blending, and perspective projection.
\end{itemize}

\section{Conclusion}

This project fully satisfies the requested academic criteria and provides a complete, integrated graphics application rather than a single-topic lab exercise. The implementation demonstrates both conceptual graphics knowledge (algorithms and transformations) and practical real-time system integration (render loop, scene architecture, input, animation, and shading).

\end{document}
